\documentclass[11pt,letterpaper]{article}
\usepackage{preamble}
\authorsig{KUKLA}
\begin{document}
\moduletitle{Q32 --- Comprehension check}{Quantifying information flow with entropy}

\begin{keybox}[Question]
A classmate shows you a plot of mutual information $I(Z;X)$ between a hidden layer
$Z$ and the input $X$ \emph{decreasing} over the course of training, and concludes
``the network is compressing away irrelevant information.'' The network is an
ordinary feed-forward net with deterministic layers.

State the single most important objection to this conclusion, and name the one
design choice in their analysis that most likely produced the downward curve.
\end{keybox}

\vspace{0.5em}
\begin{keybox}[Answer]
\textbf{Objection:} for a \emph{deterministic} map $Z=f(X)$, the mutual information
$I(Z;X)$ is not actually finite/changing in the way plotted --- it is maximal
(equal to $H(X)$, or infinite for continuous $X$). Any finite, moving number is a
property of the \emph{estimator}, not of the network.

\textbf{The design choice:} the \emph{binning / discretisation} used to estimate the
entropy of a continuous activation (together with the choice of a saturating
nonlinearity such as \texttt{tanh}). Saxe et al.\ (2018) showed this is what
manufactures the apparent ``compression phase''; switching to \texttt{ReLU} or a
finer binning typically makes it shrink or disappear.

\emph{Full-credit answer} names both: (1) deterministic $\Rightarrow$ MI is
estimator-defined, and (2) the binning/nonlinearity choice drives the curve.
\end{keybox}

\end{document}
